\fig{fv:h4:boiling}{体内气泡能否增长，取决于蒸气压能否抵住局部外压}{\begin{tikzpicture}[scale=.9]
\draw[thick](0,3.1)--(0,0)--(4.3,0)--(4.3,3.1);\fill[blue!9](0,0)rectangle(4.3,2.5);\draw[blue](0,2.5)--(4.3,2.5);
\draw[thick](2.1,1)circle(.5);\node at(2.1,1.2){$p_{\rm in}$};\draw[dim](2.1,1)--(2.6,1)node[midway,below]{$r$};\draw[->](1.1,1)--(1.6,1);\draw[->](3.1,1)--(2.6,1);
\draw[<->](4.8,1)--(4.8,2.5)node[midway,right]{$h$};\node at(2.2,3.5){$p_0$};\end{tikzpicture}}
图\ref{fv:h4:boiling}中，气泡中心深度$h$决定周围液体压强，半径$r$决定曲率附加压强。较小的泡需要更大的内外压差，不能只看液面处的大气压。若泡内含不凝结空气，同温扩张会降低其空气分压；水蒸气分压则在相平衡条件下仍由$T$决定。两种组分对体积变化的响应不同，必须分别处理。
